% !TEX program = lualatex
% ECON 201, Worksheet 6: The Cost of Making a Car. Compile with LuaLaTeX so
% the PDF is tagged (see syllabus/Econ201-Syllabus.tex for the same setup).
% Build from the course root:
%   latexmk -lualatex -cd worksheets/worksheet06.tex && latexmk -c -cd worksheets/worksheet06.tex
% Every number here is computed and printed by slides/figures/make_lecture6.py
% so the worksheet and the deck quote the same figures.
\DocumentMetadata{
  pdfstandard=UA-2,
  pdfversion=2.0,
  lang=en-US,
  tagging=on
}
\documentclass[11pt]{article}

\usepackage[letterpaper, margin=0.8in, top=0.6in, bottom=0.55in]{geometry}
\usepackage{fontspec}
\setmainfont{Fira Sans}
\usepackage{amsmath}
\usepackage{amssymb}
\usepackage{graphicx}
\usepackage{booktabs}
\usepackage{array}
% Questions are labelled by hand: under tagging=on, enumitem's enumerate
% did not step its counter, so every item printed the same label.
\newcounter{q}
\newcommand{\Q}[1]{\stepcounter{q}\par\needspace{5\baselineskip}\vspace{3pt}\noindent\hangindent=1.7em\hangafter=1
  \makebox[1.7em][l]{(\alph{q})}#1\par}
\usepackage{xcolor}
\usepackage{needspace}
\usepackage[most]{tcolorbox}
\definecolor{cream}{HTML}{FDF3EA}
\definecolor{accent}{HTML}{BF5700}
\usepackage{titlesec}
\titleformat{\section}{\large\bfseries\color{accent}}{}{0pt}{}
\titlespacing*{\section}{0pt}{6pt}{2pt}
\usepackage{hyperref}
\hypersetup{pdftitle={ECON 201 Worksheet: The Cost of Making a Car}, pdfauthor={Div Bhagia}, hidelinks}
\pagestyle{empty}
\setlength{\parindent}{0pt}
\setlength{\parskip}{3pt}

% Ruled answer space: n lines. Plain TeX loop; pgffor's \foreach breaks the
% enumerate counter when used inside an item.
\newcount\answerlinecount
\newcommand{\answerlines}[1]{%
  \par\vspace{1pt}\answerlinecount=#1
  \loop\ifnum\answerlinecount>0
    \rule{\linewidth}{0.3pt}\par\vspace{5pt}%
    \advance\answerlinecount by -1
  \repeat
  \vspace{-5pt}}

\begin{document}

{\LARGE\bfseries\color{accent} Worksheet: The Cost of Making a Car}\hfill{\large ECON 201}

% Terms sit on the cream panel, where the brand orange falls below 4.5:1;
% the darker orange keeps WCAG AA contrast (see CLAUDE.md).
\definecolor{accentdark}{HTML}{8F4100}
\newcommand{\term}[1]{\textcolor{accentdark}{\textbf{#1}}}

\begin{tcolorbox}[colback=cream, colframe=accent, boxrule=0.6pt, arc=2pt,
  left=6pt, right=6pt, top=4pt, bottom=4pt, title=Definitions,
  fonttitle=\bfseries, coltitle=white]
\term{Cost function}, $C(Q)$: the total cost of producing $Q$ units of
output. Suppose a firm has fixed costs $F$, and its variable costs are
directly proportional to the quantity it produces. Then the total cost of
producing $Q$ units is
\[
  C(Q) = F + cQ
\]
where $F$ is the \term{fixed cost}, paid regardless of what $Q$ is, and
$cQ$ is the \term{variable cost}: $c$ per unit, times $Q$ units.

\term{Average cost}, AC: the cost per unit,
\[
  \text{AC} = \frac{C(Q)}{Q}
\]
\term{Marginal cost}, MC: the increase in total cost from producing one more
unit,
\[
  \text{MC} = \frac{\Delta C}{\Delta Q}
\]
where $\Delta$ means ``the change in'': $\Delta C$ is the change in total
cost when output changes by $\Delta Q$.
\end{tcolorbox}

\needspace{10\baselineskip}\section{1. Fixed or variable?}\setcounter{q}{0}

Beautiful Cars is a small firm that manufactures specialty cars. It has to
pay for the following inputs to make its cars. Select all the costs that
change with the number of cars the firm makes each day.

\newcommand{\tick}[1]{\par\noindent\hangindent=1.7em\hangafter=1\makebox[1.7em][l]{$\square$}#1\par}
\tick{A factory equipped with machinery to assemble the cars.}
\tick{Raw materials and components.}
\tick{Production workers to operate the equipment.}
\tick{Other workers to manage the production process, and to market and sell the finished cars.}
\tick{Research and development to develop future models.}
\vspace{4pt}

\needspace{12\baselineskip}\section{2. The cost function of Beautiful Cars}\setcounter{q}{0}

Beautiful Cars has fixed costs of \$60,000 per day, and each extra car costs
\$10,000, so its cost function is
\[
  C(Q) = 60{,}000 + 10{,}000\,Q
\]
Fill in total cost, average cost, and marginal cost at each output. For
marginal cost, use the cost of one more car.

% A ruled grid rather than a booktabs table: students write in the cells,
% so every cell gets visible borders and the blank rows get extra height.
\begin{center}
\renewcommand{\arraystretch}{1.5}
\newcolumntype{C}{>{\centering\arraybackslash}p{3.0cm}}
\begin{tabular}{|>{\raggedright\arraybackslash}p{4.2cm}|C|C|C|}
\hline
Cars per day, $Q$ & 10 & 20 & 50 \\
\hline
Total cost, $C(Q)$ & \rule{0pt}{2.4em} & & \\
\hline
Average cost, $\text{AC}$ & \rule{0pt}{2.4em} & & \\
\hline
Marginal cost, $\text{MC}$ & \rule{0pt}{2.4em} & & \\
\hline
\end{tabular}
\end{center}

Average cost falls as $Q$ rises, while marginal cost does not change.
Explain why in a sentence or two.
\answerlines{3}

\end{document}
